% to be added: introductory paragraph

Hardware accelerators are widely integrated into electronic systems to offload application-specific workloads from general-purpose processors. This is especially relevant for Internet-of-Things (IoT) edge nodes, where compute and memory resources are limited, and battery lifetime is a primary constraint \cite{9985008}. Much of the work focusing on modern machine learning has been targeting convolutional neural networks (CNNs) \cite{Lecun2015436}, which provide state-of-the-art accuracy. However, this performance comes at a price, as neural networks typically require large amounts of data and computational resources \cite{GranmoOle-Christoffer2018TTM-}, which is a difficult match for the resource constrained nature of edge devices. 


% pagraph presenting the Tsetlin Machine, quick introduction of the architecture and how it works, and why it is interesting for hardware acceleration and edge deployment. This will be expanded in the theory chapter, but a quick introduction here will help set the stage for the rest of the thesis.
The Tsetlin Machine (TM) offers a qualitatively different approach to other machine learning models. Introduced by Ole-Christoffer Granmo in 2018 \cite{GranmoOle-Christoffer2018TTM-}, the TM is a logic-based learning algorithm founded on propositional logic and finite-state automata. It constructs propositional formulas (clauses) to perform classification tasks, where each clause is a conjunction of literals (input features and their negations) that vote for or against a particular class. The learning process involves updating the states of Tsetlin Automata based on feedback from the classification results, allowing the model to dynamically refine its decision boundaries. The architecture of the TM enables execution using simple bitwise operations, making it particularly suitable for hardware acceleration and edge deployment \cite{AndersTunheimSvein2025AA86, TunheimSveinAnders2023CTMT, TunheimSveinAnders2025TMIC}. Fixed-function FPGA and ASIC accelerators have demonstrated that TM inference can be executed at high throughput and low energy \cite{AndersTunheimSvein2025AA86, TunheimSveinAnders2023CTMT, TunheimSveinAnders2025TMIC}, but these designs commit to a specific TM architecture at synthesis time and require significant redesign to accommodate new variants or problem configurations.

Processor-oriented acceleration preserves programmability by embedding domain-specific logic inside a general-purpose processor rather than alongside it. The RISC-V Instruction Set Architecture (ISA) is well-suited to this approach, providing reserved opcode space for user-defined instruction extensions that integrate directly into the processor pipeline \cite{WatermanLee2011,riscv_unpriv_isa}. A preceding project thesis \cite{JB2025} demonstrated this on the Noisy XOR (Section \ref{sec:other_datasets}) problem, implementing a custom clause evaluation instruction on the NEORV32 soft-core processor and achieving a $4.42\times$ inference speedup at a cost of only 18 additional LUTs. 

This thesis builds on that work by examining how the approach scales to a realistic workload, specifically MNIST digit classification. Whereas the previous work operated on a model of 480 bytes, this work targets a model of approximately 3.81\,MB, an increase of roughly $8000\times$. A scaling of this magnitude introduces a qualitatively different challenge: the model far exceeds the processor's internal memories, placing memory bandwidth rather than computation at the centre of the performance problem.


\section{UN Sustainability Goals}
This work relates to two of the United Nations Sustainable Development Goals~\cite{UNSDG}. The first is Goal 7, Affordable and Clean Energy. As machine learning inference on edge devices becomes more widespread, the energy it consumes is a growing concern, particularly for battery-powered IoT nodes. The Tsetlin Machine performs inference using only bitwise operations and integer arithmetic, avoiding the multiply-accumulate and floating-point operations that dominate the energy cost of conventional neural networks. Investigating how such a model can be executed efficiently on low-power hardware contributes to reducing the energy demand of on-device intelligence. The second is Goal 9, Industry, Innovation and Infrastructure. This work is built on open and royalty-free foundations, namely the RISC-V instruction set architecture and the open-source NEORV32 soft-core processor, which lower the barrier to developing custom computing hardware. Exploring how application-specific acceleration can be added to an open processor rather than a proprietary platform supports innovation in accessible and adaptable computing
infrastructure.


\section{Related Work}  
\label{sec:related_work}

\subsection{Hardware Accelerators for Tsetlin Machines}
\label{sec:related_hw}

The first realized HW implementation of the TM was introduced in 2020. It was implemented both on 65 nm ASIC as well as on an Altera Stratix IV FPGA \cite{WheeldonAdrian2020Labe}. The design supported on-device learning by utilizing parallel clause updates and stochastic feedback, by using lightweight pseudo-random number generators. It demonstrated competitive energy efficiency, outperforming a Convolutional Binary Neural Network by $2.5\times$ in energy per operation. However, the accuracy fell short of the software equivalent of the TM. This difference was attributed to differences in random number generator precision.

Extending this to spatial data, the first hardware accelerator for the Convolutional TM (CTM) \cite{GranmoOle-Christoffer2019TCTM} was presented in 2023 \cite{TunheimSveinAnders2023CTMT}. The design targeted binary image classification and supported both inference and on-device learning, and evaluted on the NoisyXOR dataset (Section \ref{sec:other_datasets}). It achieved a throughput 13.3 times higher than a software implementation running on a 2.7 GHz CPU. The work also identified the quality of the random number generator as a critical factor in achieving accuracy comparable to software, finding a 16-bit LFSR to be the optimal configuration. 

More recently, two hardware accelerators for the Convolutional Coalesced TM (ConvCoTM) \cite{GlimsdalSondre2021CMTM} were presented in 2025. The first is an FPGA implementation targeting $28 \times 28$ grayscale images across ten output classes, achieving 134k classifications per second and $13.3\,\mu$J energy per classification (EPC) on a Xilinx Zynq ZU7EV while supporting full on-device learning \cite{TunheimSveinAnders2025TMIC}. The design was evaluated on MNIST, Fashion-MNIST and Kuzushiji-MNIST (Section~\ref{sec:other_datasets}), maintaining competitive accuracy across all three datasets.The second is a 65 nm ASIC inference-only implementation of the same architecture, achieving 60.3k classifications per second at just 8.6 nJ EPC, one of the lowest reported for TM hardware \cite{AndersTunheimSvein2025AA86}. Together, these designs demonstrate how dedicated TM hardware can deliver high energy efficiency and competitive accuracy for edge deployment. 

\subsection{Processor-Oriented TM Acceleration}
\label{sec:related_processor}

Zeng et al.\ propose an optimized software TM inference implementation targeting general-purpose CPUs, evaluated using the gem5 simulator with an ARM processor architecture \cite{zeng2025fastcompacttsetlinmachine}. The key contributions are a bitwise operation based clause evaluation method, an early exit mechanism that terminates clause evaluation as soon as a failing literal group is detected, and a Reorder strategy that reorders literals and TA actions before inference to maximize the likelihood of early exits. Together, these techniques reduce inference time by up to 96.71\% compared to a conventional integer-based baseline, with the most significant gains observed on high-dimensional datasets such as MNIST (Section \ref{sec:mnist}). Because the approach is purely software-based, it requires no hardware modifications and is directly portable across processor platforms.

\subsection{RISC-V Custom Extensions for Machine Learning}
\label{sec:related_riscv}

Wang et al.\ propose a set of custom RISC-V SIMD instructions targeting the convolution, activation, and pooling operations common in neural network inference \cite{WangShihang2024OCCU}. Rather than relying on a large dedicated accelerator, the approach adds a compact set of specialized instructions that process multiple data values in parallel, achieving meaningful speed and energy efficiency improvements on FPGA-based edge systems while keeping the processor design lightweight and programmable. This work is directly relevant to the approach taken in this thesis, where a custom RISC-V instruction is used to accelerate TM clause evaluation within the NEORV32 processor (Section \ref{sec:neorv32}), targeting a similar balance between programmability and hardware efficiency.

\subsection{Prior Work: Processor-Level Acceleration of Tsetlin Machines
Using Custom RISC-V ISA Extensions}
\label{sec:related_prior}

The present thesis extends the work reported in \cite{JB2025}, which investigated processor-level TM inference acceleration on the same NEORV32 RISC-V processor deployed on a Xilinx Artix-7 FPGA (Section \ref{sec:neorv32}). That work designed a custom clause evaluation instruction using the CFU, implemented using the R3-type instruction format available in the version of NEORV32 used at the time (Section \ref{sec:neorv32_cfu}), and evaluated inference performance on the NoisyXOR dataset.

The results demonstrated a clear speedup for the NoisyXOR problem, where the model is small enough to reside entirely in on-chip DMEM, confirming that the CFU-based approach can meaningfully reduce clause evaluation time when memory access is not the bottleneck. The central question left open by that work is whether the same acceleration benefit scales to larger, more practically relevant problems such as MNIST digit classification, where the model size far exceeds on-chip memory capacity and must be stored in off-chip DDR2 SDRAM (Section \ref{sec:ddr2}). Investigating this question is the primary research objective of the present thesis.